\section{Business Cycles, Innovation, Crisis}

\begin{frame}{Business Cycles}
    \mytext{
    \showmytikz[1,2,3]{../shared/tikz/tikz-porkcycle}
    
    Economic growth is not steady, we see phases of expansion and contraction, of boom and bust.

    What exactly determines these cycles? Are they predictable?
    }
    \mycomment{
    In stock broker jargon, markets with rising prices (expanding markets) are nicknamed \enquote{bull markets} -- the bull is a fierce animal \enquote{fighting upwards}. Falling prices (contracting markets) are nicknamed \enquote{bear markets} -- the bear is a fierce animal \enquote{fighting downwards}.

    Longing: you \enquote{long} something (you haven't got it, but you want to have it) -- you are betting on rising prices, you are trying to buy. Long selling: you actually own the asset that you are selling. 

    Shorting: you are trying to sell. Short selling: you are trying to sell without acutally having bought the asset yet -- you should have the option to buy it though (call option). 
    }
\end{frame}

\begin{frame}{Anticyclic Policy}
    \mycomment{
    If you know how the cycle runs, you could (in theory) counter-act it, thus evening the amplitude out. If ups and downs are natural and deterministic for lassaiz-faire supply and demand, adding artificial government demand (and supply) could smooth these out. This is the idea of anticyclic fiscal policy.
    }
        \showmytikz[1,2,3]{../shared/tikz/tikz-anticyclic.tex}

\end{frame}

\begin{frame}{Kitchin Cycle}
    \mytext{
    \showmytikz[1]{../shared/tikz/tikz-businesscycles}
    
    The Kitchin cycle is an inventory cycle, like the Pork Cycle, with a cycle length of 3-5 years.
    }
\end{frame}

\begin{frame}{Juglar Cycle}
    \mytext{
    \showmytikz[2]{../shared/tikz/tikz-businesscycles}
    
    The Juglar cycle is a fixed investment cycle. 

    Cycle length: 7-11 years.
    }

\end{frame}

\begin{frame}{Overlapping Cycles}
    \mytext{
    \showmytikz[1,2,3]{../shared/tikz/tikz-businesscycles}
    
    The red curve shows the sum of the cycles already introduced (with a scaled down amplitude).  
    }
\end{frame}

\begin{frame}{Kuznets Swing}
    \mytext{
    \showmytikz[4,5]{../shared/tikz/tikz-businesscycles}
    
    Kuznets swings relates to infrastructure investment.

    Cycle length: 15-25 years.
    }
\end{frame}

\begin{frame}{Kondratiev Wave}
    \mycomment{
    Kondratiev Waves are larger boom and bust cycles attributed to major technological breakthroughs. The invention of the steam engine, the combustion engine, the computer are much quoted.

    In these discussions, the negative (contractive) phase is often left out. But every new invention also makes existing investment obsolete, causing factories to shut etc. and thus having a negative impact as well.
    }
    \mytext{
    \showmytikz[6,7]{../shared/tikz/tikz-businesscycles}
    
    Kondratiev Waves are (legendary?) long-term ups and downs due to technological change. 

    Cycle length: 45-60 years.
    }
\end{frame}

\begin{frame}{Business Cylces: Random Walk}
    \mytext{
    \showmytikz[7,8]{../shared/tikz/tikz-businesscycles}
    In reality, economic activity appears to resemble a random walk (here in green) rather than cycles. If that is so, there is no more argument for a countercyclical policy. 
    }
    \mycomment{
    There is a real difference between \enquote{ordinary} business cycles, and black-swan crisis moments. Random ups and downs might rock the boat. But every now and then, for whatever reasons, waves may sum up creating a perfect storm and capsizing a boat. 

    The Greek word \enquote{crisis}  means the \textit{turning point} of an illness, where the patient either recovers or not.  The illness itself is not a crisis in the strict sense of the word.
    }
\end{frame}

\begin{frame}{Business Cycles: Credit Driven}
    \mytext{
    Business cycles could be the result of credit boom and bust, and of the central bank's interest setting behaviour (known as Austrian Business Cycle Theory).

    Then it would be the government \enquote{rocking the boat}, depressing interest rates below their natural level, thus setting wrong price signals (over-optimistic investment, \enquote{malinvestment}). 

    The crisis can then not be fixed by further expansion, but is a necessary correction.
    }
\end{frame}


\begin{frame}{Creative Descruction}
    \showmytikz{../shared/tikz/tikz-schumpeter}
    \mycomment{
    The new technology, shown in green, overtakes the old technology, shown in red. Total output increases seamlessly since the new technology is more productive than the old one.
    }
    
\end{frame}